\begin{statement}{A.5.18}{Proved}{Every horizontal tacnode image is elliptic}
Let $f:Y\to C_0$ be a proper holomorphic map from a smooth
compact complex threefold with $b_2(Y)=0$, with connected
fibres and smooth general fibre a complex two-torus.
Let $mS$ be a full fibre with irreducible reduced support,
$\eta:T\to S$ its normalization, and $\pi:Z\to T$ a
resolution with an irrational ruling $p:Z\to B$.
Suppose every exceptional curve of $\pi$ is vertical and
the conductor divisor on $Z$ is $2C+V$, where $C$ is a
section of $p$ and $V$ is vertical. Let $D=\eta\pi(C)$.
Suppose the generic plane-curve transversal to $D$ in $S$ is an
ordinary tacnode.

Then $D$ is a smooth elliptic curve, $B=C$ is elliptic,
and $C\to D$ is an unramified double cover. The transversal
type along $D$ is everywhere the same ordinary tacnode.
The normalization $T$ is a smooth minimal ruled surface
over this elliptic curve. Its conductor is twice a
section, and $e(S)=0$. In particular its minimal
resolution is the identity.
In particular a rational image of
the horizontal tacnodal conductor cannot occur. The proof
allows vertical conductor and rational singularities of
$T$; it does not assume that the normalization map is
immersive at the special points under examination.
\end{statement}
\begin{proof}
If the normalization of $D$ has positive genus, A.5.14
and A.5.5 make $D$ smooth elliptic and its transversal
type constant. There are two generic normalization
branches of a tacnode. A vertical component of $Z$
is rational, so it cannot dominate this elliptic curve.
Thus both branches lie on $C$, and $C\to D$ is an
unramified double cover: the constant local equation
$x^2-y^4=0$ has two separately smooth normalization
branches on every sufficiently small disc in $D$.
Riemann--Hurwitz gives $g(C)=g(B)=1$.
It remains to exclude the case $D^\nu=\Sigma\simeq\PP^1$.

The finite map $C\to\Sigma$ has degree at least two,
because $C\simeq B$ has positive genus. Since a generic
tacnode has exactly two normalization branches, this
degree is exactly two. Consequently no vertical
conductor component can also dominate $D$: its generic
points would supply an additional normalization branch.
Write $g=g(C)\ge1$ and let $\mathcal B\subset\Sigma$
be the reduced branch divisor of this double cover.
Riemann--Hurwitz gives
\[
                   \deg\mathcal B=2g+2.             \tag{A.T.1}
\]

We first bound the number of branches of the singular
image curve itself. A.5.8 supplies the descended ruling
$\widetilde p:T\to B$. Identify $C$ with $B$ using the
section, let $\phi:B\to D$ be its actual map to the
image curve, and put
\[
                    \psi=\phi\widetilde p:T\to D.
\]
Every vertical conductor component has constant $\psi$
value. Its image under $\eta$ differs from $D$, as just
proved, and meets $D$ in finitely many points. Thus
there is a finite set $A\subset D$ such that, for each
$s\in S$, either $\psi(\eta^{-1}(s))$ is a singleton
or all its values belong to $A$. Include in $A$ the
constant values on the vertical conductor components
and the values of $\phi$ over their finitely many
intersections with $D$. Fibres over the remaining
points of $D$ have their preimages on $C$, where all
the $\psi$ values equal that same point of $D$.

Choose an embedded finite tree $J\subset D$ containing
$A$, in a triangulation of the compact analytic curve.
It can be constructed by joining these finitely many
points by arcs, subdividing their intersections and
retaining a spanning tree. The quotient
$q:D\to D/J$ is a homotopy equivalence: a contractible
CW subcomplex is a cofibration, and the homotopy
extension property extends its contraction. Since
$q\psi$ is constant on every fibre of the finite
proper quotient $\eta$, there is a continuous map
\[
 a:S\to D/J,\qquad a\eta=q\psi,\qquad a|_D=q.     \tag{A.T.2}
\]
The last equality holds first on the dense regular
locus of $D$ and then everywhere by continuity.

Use the explicit triangulated neighbourhood retraction
$r:N\to S$ of (A.Pi.3). For a sufficiently small disc
$\Delta$ around the fibre value, properness gives
$U=f^{-1}(\Delta)\subset N$. The map $(ar)|_U$ is
surjective on fundamental groups because its restriction
to $D$ is the homotopy equivalence $q$. Perturbing loops
off the real-codimension-two analytic support gives
$\pi_1(U\setminus S)\twoheadrightarrow\pi_1(U)$.
The torus-bundle exact sequence over $\Delta^*$ therefore
gives the actual quotient maps
\[
 \Z^4\rtimes\Z\twoheadrightarrow\pi_1(U)
       \twoheadrightarrow\pi_1(D/J)\simeq\pi_1(D).
                                                        \tag{A.T.3}
\]
The first group is metabelian, so the last is metabelian.

For each singular point $x$ of $D$, let $r_x$ be the
number of its normalization branches and put
$\beta=\sum_x(r_x-1)$. The finite normalization map
identifies exactly those finite sets of points on
$\PP^1$; since it is a compact Hausdorff quotient,
this is also the actual underlying topological quotient.
Adjoining arcs between the points to be identified
gives a CW presentation homotopy equivalent to
$S^2\vee\bigvee^\beta S^1$. Hence
$\pi_1(D)$ is the free group $F_\beta$.
For $\beta\ge2$ this group surjects onto $S_4$ by
sending two free generators to $(1234)$ and $(12)$;
the commutator $A_4$ of that quotient is nonabelian.
Equation (A.T.3) therefore forces $\beta\le1$.
By A.5.14 the total delta invariant $\delta=p_a(D)$
is at most two. If $s$ denotes the number of points
of $\Sigma$ over $\operatorname{Sing}D$, then
\[
 s=\#\operatorname{Sing}D+\beta\le\delta+\beta\le3.
                                                        \tag{A.T.4}
\]

First suppose that $D$ is a local complete intersection.
We retain the exact bundles and jets of A.5.5.
Let $\nu:\Sigma\to D$ be normalization, $I=\mathcal I_D$,
\[
 N=(\nu^*(I/I^2))^*,\qquad
 M=\nu^*(\mathcal O_Y(S)|_D).
\]
Thus $M$ is a degree-zero line bundle on $\Sigma$.
Let $\Xi$ be the curve-normalization
conductor divisor, $c=\deg\Xi=2\delta$, and let
$R\subset N$ and $Q=N/R$ be the saturated radical and
quotient of the generic rank-one Hessian. Write
$q_0=\deg Q$. The adjunction and conductor comparisons
(A.J.1) give
\[
 \deg N=-2+c,\quad \deg R=-2+c-q_0,\quad
 0\ne a\in H^0(\Sigma,M\otimes Q^{-2}),\quad q_0\le0.
                                                        \tag{A.T.5}
\]
Here $a$ denotes the intrinsic quadratic coefficient,
whose zero divisor has degree $-2q_0$.
The actual quartic invariant $\Delta=ae-b^2/4$ of
A.5.5 is meromorphic in
$L^2$, where $L=M\otimes Q^{-1}\otimes R^{-2}$.
Its pole bound (A.J.9) gives a nonzero holomorphic
section $\sigma$ of the actual square line bundle
$L^2(4\Xi)=(L(2\Xi))^2$. Its zero divisor has degree
\[
               d_\Delta=8+2q_0.                    \tag{A.T.6}
\]

The double cover $C\to\Sigma$ is obtained by adjoining
$\sqrt{-\Delta}$ to the function field, with the line
trivializations specified by $L$. Indeed over the
generic point a tacnode has normal expansion
$a\ell^2+b\ell r^2+er^4$, and the two normalization
branches are distinguished by the two roots of
$a(\ell/r^2)^2+b(\ell/r^2)+e$. Their discriminant
is $-4\Delta$. Changing a local frame of $L$ multiplies
the representing function of $\Delta$ by a square,
so this describes the same degree-two extension on
every overlap. At a point with local parameter $z$,
write that representing function as $z^n u(z)$,
$u(0)\ne0$. A root of the unit and removal of the
even power of $z$ show that its normalization is
ramified exactly when $n$ is odd. Adding $4\Xi$ changes
each order by an even number. Thus $\mathcal B$ is
exactly the odd-multiplicity support of
$\operatorname{div}(\sigma)$, and
\[
                  2g+2\le 8+2q_0.                 \tag{A.T.7}
\]

At a point of $D_{\rm reg}$ where $a\ne0$, A.5.17
applies to the actual local $m$th root of the full-fibre
map. Its nonzero Hessian makes the local equation a
constant transverse product, necessarily of type
$A_3$ because this is the generic type along $D$.
The conductor cover is therefore unramified there.
Every branch point outside $\nu^{-1}(\operatorname{Sing}D)$
is consequently a zero of $a$. Equations (A.T.4)--(A.T.5)
give the second bound
\[
                    2g+2\le -2q_0+s.              \tag{A.T.8}
\]
Adding (A.T.7) and (A.T.8) yields
$4g+4\le8+s\le11$. Since $g\ge1$, we have
\[
           g=1,\qquad \deg\mathcal B=4,
                    q_0\in\{-2,-1,0\}.            \tag{A.T.9}
\]

We next prove the exact local obstruction required in
the equality cases. Let a smooth point of $D$ be a
branch point at which $\Delta$ has a simple zero.
The preceding nonzero-Hessian argument gives $a(0)=0$.
In coordinates $D=(x,y)=0$ with curve parameter $z$,
choose the holomorphic saturated Hessian factor as the
first normal coordinate. The quadratic jet is then
$a(z)x^2$. The radical cubic vanishes identically,
so a local reduced equation has the exact expansion
\[
 f_0=a(z)x^2+x^2(xA+yB)+xy^2b(z)+xy^3C+y^4E,
                                                        \tag{A.T.10}
\]
with holomorphic coefficients. The identity
$\Delta=ae-b^2/4$, $a(0)=0$ and the simple zero of
$\Delta$ imply $b(0)=0$, $a'(0)\ne0$ and $e(0)\ne0$.
Indeed an order at least two for $a$, together with
$b(0)=0$, would give an order at least two for $\Delta$;
and then $\Delta'(0)=a'(0)e(0)$ is nonzero.
Thus $b/a$ is holomorphic and $E(0)\ne0$.

Set $X=x+b(z)y^2/(2a(z))$. In this actual analytic
coordinate change the coefficient of $Xy^2$ cancels.
All terms created from $x^2(xA+yB)$ lie in
$(X^2)+(Xy^3)+(y^4)$. Therefore, collecting the
terms of $X$-degree zero and one, gives exactly
\[
 f_0=X^2A_1(X,y,z)+Xy^3B_1(y,z)+y^4C_1(y,z),
 \qquad C_1(0)\ne0.
\]
Choose a fourth root of $C_1$ and define
\[
 Y=yC_1^{1/4}+
                  \frac{XB_1}{4C_1^{3/4}}.
                                                        \tag{A.T.11}
\]
The map $(X,y,z)\mapsto(X,Y,z)$ is locally
biholomorphic, and direct expansion gives
$f_0-Y^4\in(X^2)$. Hence $f_0=Y^4+X^2H(X,Y,z)$.
Its quadratic jet along $D$ is still $a(z)X^2$,
so $H(0,0,z)=a(z)$ and $\partial H/\partial z(0)\ne0$.
The further coordinate $Z=H$ gives exactly
\[
                         f_0=Y^4+X^2Z.             \tag{A.T.12}
\]
To apply this to the actual map, first take a local
$m$th root of its unit. This multiplies the reduced
equation by a unit, which changes neither zero order
of $a$ nor of $\Delta$. The same coordinate construction
then makes the full map $(Y^4+X^2Z)^m$.
Statement A.5.15 excludes it. Consequently every
smooth branch point has order at least three in
$\operatorname{div}(\sigma)$.

We also record precisely when the conductor denominator
can be removed at a singular point $x\in D$.
If the intrinsic section $a$ is nonzero at every
normalization point above $x$, choose an actual
constant linear combination of the two lci conormal
generators which is nonzero under the normalized
quotient at all those finitely many points. The
corresponding diagonal entry $a_0$ of the symmetric
Hessian matrix belongs to $\mathcal O_{D,x}$ and
has nonzero value on every normalization branch.
It is therefore a unit in $\mathcal O_{D,x}$.
The other normalized factor coefficient is the
actual quotient $h=q_{XY}/a_0\in\mathcal O_{D,x}$.
Thus square completion and monic polynomial division
in the proof of (A.J.7)--(A.J.8) both take place
over $\mathcal O_{D,x}$ itself. The regular-sequence
graded pieces inject into their normalization
pullbacks, so their cancellations already descend
on $D$. We may repeat that proof with $d=1$.
It gives a holomorphic section $\Delta$ of $L^2$
above $x$, without a conductor pole. At a branch
point above $x$, this proves the sharper inequality
\[
 \operatorname{ord}\sigma\ge
                        4\operatorname{ord}\Xi+1.
                                                        \tag{A.T.13}
\]

We now exhaust the three integer cases in (A.T.9).
If $q_0=-2$, the divisor of $\sigma$ has degree four
and has four odd-multiplicity points, so each is
simple. At most three lie above singularities of
$D$, by (A.T.4). The remaining smooth branch point
contradicts (A.T.12).
If $q_0=0$, the section $a$ has no zeros and all
four branch points would have to lie above the
singular locus; this contradicts $s\le3$.

Finally let $q_0=-1$, so that $\deg\operatorname{div}(a)=2$
and $d_\Delta=6$. If $k$ of the four branch points
lie above $\operatorname{Sing}D$, then the local
obstruction and parity give
\[
               6\ge k+3(4-k)=12-2k.
\]
Hence $k=s=3$, and equality holds: the three singular
preimages have order one in $\operatorname{div}(\sigma)$
and the one smooth branch point has order three.
The relations $s\le\delta+\beta\le3$ now force
$\delta=2$, $\beta=1$ and two singular points of
delta invariant one. The exact delta-one
classification in A.5.14 identifies them as one
ordinary node and one cusp. Their normalization
conductor orders are respectively $(1,1)$ and $2$;
these follow directly from the displayed node and
cusp rings in that proof.

At the cusp preimage, if $a$ were nonzero,
(A.T.13) would make the order of $\sigma$ at least
$4\cdot2+1=9$, exceeding its total degree six.
So $a$ vanishes there. It also vanishes at the
smooth branch point. These two distinct points
exhaust the degree-two divisor of $a$. The section
is consequently nonzero at both node preimages.
Apply (A.T.13) at the node: each of its two branch
points has order at least $4\cdot1+1=5$ in
$\operatorname{div}(\sigma)$, again exceeding the
total degree six. This excludes every rational lci image.

We finish by treating all non-lci image germs permitted
by $\delta\le2$ and $\beta\le1$, with explicit analytic
rings and transverse coordinate maps. For a unibranch
reduced curve ring $A\subset\C\{t\}$ of colength at
most two, an element of order one makes $A=\C\{t\}$.
Otherwise its least positive valuation is two or three:
a least valuation at least four leaves three linearly
independent missing coefficients. If it is three and
the colength is two, the inclusion
$A\subset\C+t^3\C\{t\}$ is equality, since the latter
already has colength two. If it is two, choose a
parameter with $t^2\in A$ by taking a square root of
the leading unit. The even part of every convergent
series belongs to $\C\{t^2\}\subset A$. Subtracting
it shows that the least odd element can be written
$t^{2k+1}v(t^2)$; division by its unit puts
$t^{2k+1}$ in $A$. Consequently
$A=\C\{t^2\}\oplus t^{2k+1}\C\{t^2\}$, of colength
$k\le2$. These are the planar cusp and the planar
$A_4$ curve, besides the smooth case.

For two branches, write their image rings as $A_1,A_2$.
The two projections of $A\subset A_1\oplus A_2$ are
surjective. Their kernels define a common finite
quotient $Q_0$ and an exact fibre-product description
$A=A_1\times_{Q_0}A_2$. Thus
$\delta=\delta_1+\delta_2+\operatorname{length}Q_0$.
If both branches are smooth, the quotient is
$\C\{t\}/(t^k)$ with $k\le2$; after changing the
second parameter, the two maps identify the same
$k$-jet. The map
\[
 x\longmapsto(t,s),\qquad
 y\longmapsto(t^k,-s^k)
\]
identifies this fibre product with
$\C\{x,y\}/(y^2-x^{2k})$: write a matching pair as
the sum of its average and $y$ times its difference
divided by $2x^k$. These are the node and tacnode.
If one branch is singular, the delta budget forces
one cusp, one smooth branch, and $Q_0=\C$.
Three branches are excluded by $\beta\le1$.

An abstract plane curve ring in the given smooth
threefold is lci by the actual transverse-coordinate
construction in A.5.14. Thus the only remaining
non-lci analytic rings, necessarily consuming the
whole delta budget at a single point, are
\[
\begin{array}{ll}
 A=\C\{t^3,t^4,t^5\},
 &I=(y^2-xz,\ yz-x^3,\ z^2-x^2y),\\[2pt]
 A=\C\{t^2,t^3\}\times_{\C}\C\{s\},
 &I=(xz,\ yz,\ y^2-x^3).
\end{array}                                             \tag{A.T.14}
\]
These are normal forms of the actual embeddings.
Their three displayed coordinate functions form a
basis of the cotangent space of each ring. Lift them
to $\mathcal O_{Y,x}$; their differentials are linearly
independent because the embedding dimension is three.
They are therefore ambient coordinates by the inverse
function theorem. The displayed ideals are their
kernels, obtained respectively by the parametrizations
$(t^3,t^4,t^5)$ and the union of $(t^2,t^3,0)$ with
$(0,0,s)$.

We show that the local reduced equation $f_0$ actually
lies in $I^2$ in each case. In the first case put
\[
 g_1=y^2-xz,\quad g_2=yz-x^3,\quad
 g_3=z^2-x^2y,\quad
 F=x^5+xy^3+z^3-3x^2yz.
\]
The exact identities
\[
 xF=g_2^2-g_1g_3,\qquad
 yF=xg_1^2+g_2g_3,\qquad
 zF=g_3^2+xg_1g_2                              \tag{A.T.15}
\]
show that the maximal ideal annihilates the class of
$F$ modulo $I^2$. Moreover
$I^{(2)}=I^2+(F)$ in the convergent local ring.
Here is a finite-module verification of that equality.
Let $J=I^2+(F)$. The quotient $\C\{x,y,z\}/J$ is
finite over the discrete valuation ring $\C\{x\}$,
since its reduction is the finite parametrized curve.
Its generic rank is nine: the curve has degree three
over the $x$ coordinate, and the square of its smooth
codimension-two ideal has length three along each
generic normal slice. Modulo $x$, the quotient has
the nine-element basis
\[
 1,y,y^2,y^3,\ z,yz,y^2z,\ z^2,yz^2,
\]
because its ideal is $(y^4,y^3z,y^2z^2,z^3)$.
A finite module over a discrete valuation ring is
the sum of its free and torsion parts; equality of
generic rank and residue-space dimension here makes
the torsion part zero. Thus $J$ is $x$-saturated.
By (A.T.15), $J$ and $I^2$ agree after inverting $x$.
The curve is lci away from the origin, so its symbolic
square differs from its ordinary square only by a
module supported there. Hence
$I^{(2)}=I^2:x^\infty=J$.

On its normalization let
$a_1=dg_1$ and $b_1=dg_2+t\,dg_1$.
They form a basis of the normalized conormal image;
$dg_3=-2t^2a_1+t b_1$.
The quadratic term of $F$ in this basis is
\[
 q_F=\frac3t a_1^2-\frac3{t^2}a_1b_1
                           +\frac1{t^3}b_1^2,
 \qquad \det(q_F)=\frac3{4t^4}.                  \tag{A.T.16}
\]
Every equation in $I^2$ has a holomorphic quadratic
form in this same basis. If $f_0=cF+I^2$ had
$c(0)\ne0$, the determinant of its quadratic form
would have the noncancellable leading pole
$3c(0)^2/(4t^4)$: the cross terms with a holomorphic
form have pole at most three. It could not have
generic rank one. Thus $c(0)=0$, and (A.T.15) gives
$f_0\in I^2$.

For the second ring put $G=y^2-x^3$ and
$P=(z,G)$, $Q_1=(x,y)$. Then $I=P\cap Q_1$ and
\[
 I^{(2)}=P^2\cap Q_1^2=I^2+(zG).                \tag{A.T.17}
\]
Indeed $P^2=(z^2,zG,G^2)$, while $zG,G^2$ already
belong to $Q_1^2$. Multiplication by $z^2$ is
injective modulo $Q_1^2$, so membership of
$z^2A+zGB+G^2C$ in $Q_1^2$ forces $A\in Q_1^2$.
This gives the asserted generating ideal directly.
Furthermore
$x(zG)=(xz)G$, $y(zG)=(yz)G$, and
$z(zG)=(yz)^2-x(xz)^2$ all lie in $I^2$.
On the cusp branch use $a_1=d(xz)=t^2dz$ and
$b_1=dG=-3t^4dx+2t^3dy$.
The quadratic form of $zG$ is $t^{-2}a_1b_1$,
of determinant $-1/(4t^4)$. The same rank-one
argument forces its coefficient to vanish at the
origin, and the three identities again give
$f_0\in I^2$.

For clarity, the bundle used at these non-lci points
is now defined by its actual morphism. Put
\[
 K=\operatorname{im}\bigl(\nu^*(I/I^2)
                         \longrightarrow\nu^*\Omega_Y^1\bigr),
 \qquad N=K^*.
\]
The coherent image $K$ is locally free of rank two
on the smooth curve $\Sigma$. Let $R_{\rm ram}$ be
the zero divisor of $d\nu$, and let $K_{\rm sat}$
be the kernel of the surjection
$\nu^*\Omega_Y^1\to\Omega_\Sigma^1(-R_{\rm ram})$.
The torsion quotient $K_{\rm sat}/K$ defines a divisor
$J_{\rm jac}$. Taking determinants in these actual
maps gives
\[
 \deg N=-2+\deg(J_{\rm jac}-R_{\rm ram}).         \tag{A.T.18}
\]
For $(t^3,t^4,t^5)$, the wedge $a_1\wedge b_1$
has coefficients $t^8(3,4t,5t^2)$, while $d\nu$
has common factor $t^2$. Thus the correction is
$8-2=6$. For the cusp branch in the second case,
the wedge is $-t^5(2,3t,0)$ and the tangent common
factor is $t$, giving correction four. On its
smooth branch the wedge is $s^2(dx\wedge dy)$,
giving correction two. In both cases (A.T.18)
therefore gives $\deg N=4$.

Because $f_0\in I^2$, its quadratic form is a
holomorphic section of $\operatorname{Sym}^2K\otimes M$,
independent of its expression as a sum of products
of ideal generators. This follows either from its
actual Hessian or from injectivity of
$\operatorname{Sym}^2K\to\nu^*\operatorname{Sym}^2\Omega_Y^1$.
Its saturated radical and quotient again give line
bundles $R,Q$, with $q_0=\deg Q\le0$ and
$\deg R=4-q_0$. The generic quartic invariant
is meromorphic in $L^2$, where
$L=M\otimes Q^{-1}\otimes R^{-2}$; now
\[
                        \deg L^2=-16+2q_0.        \tag{A.T.19}
\]

We bound its poles using actual transverse affine
coordinate maps, with the curve parameter held
fixed during every normal differentiation. For the
space cusp use, for $t\ne0$,
\[
 (u,v,t)\longmapsto
 \left(t^3,\ t^4+\frac{v}{3t^5},\
                   t^5-\frac{u}{t^3}+\frac{2v}{3t^4}\right).
                                                        \tag{A.T.20}
\]
This is locally biholomorphic and its two transverse
derivatives are the dual frame to $a_1,b_1$.
Each coefficient of a derivative has pole at most
five. The first normal derivatives of all three
$g_i$ are holomorphic, since
$dg_1=a_1$, $dg_2=b_1-ta_1$ and
$dg_3=-2t^2a_1+t b_1$.
Second normal derivatives have pole at most ten.
For $f_0=\sum c_{ij}g_i g_j$, the cubic Taylor
coefficients consequently have pole at most ten:
the only surviving terms use a first and a second
derivative of the $g_i$, or a first derivative of
$c_{ij}$ and two first derivatives of the $g_i$.
Every fourth derivative of a holomorphic ambient
function in (A.T.20) has pole at most twenty.

Choose a holomorphic change of the $N$ frame adapted
to its saturated radical and quotient. It does not
increase these bounds. The radical cubic is zero
because the generic transversal is $A_3$.
Writing the relevant Taylor terms as
$a\ell^2+b\ell r^2+er^4$ therefore gives a
holomorphic $a$, a mixed cubic $b$ of pole at most
ten, and $e$ of pole at most twenty. The invariant
$\Delta=ae-b^2/4$ has pole at most twenty.
These are the same intrinsic jets on overlaps:
to weighted order four, where $\ell$ has weight
two and $r$ has weight one, an adapted change is
$\ell'=\alpha\ell+u r^2$, $r'=\beta r$.
Direct substitution multiplies $\Delta$ by
$\alpha^{-2}\beta^{-4}$; changes of the defining
section supply the factor $M^2$. Higher weighted
terms and changes of the curve parameter do not
alter these coefficients. This verifies the
transition law for the line $L^2$ in (A.T.19).

There is only one singular preimage $p$ in this case.
Thus $\Delta$ is a holomorphic section of the square
line $L^2(20p)$, whose degree is $4+2q_0$.
As in (A.T.7), the even pole correction preserves
branch parity, so
\[
 \deg\mathcal B\le4+2q_0,
 \qquad \deg\mathcal B\le-2q_0+1.
\]
The second bound again uses A.5.17 away from that
single singular image preimage. Adding them gives
$2\deg\mathcal B\le5$, contradicting
$\deg\mathcal B=2g+2\ge4$.

For the cusp-plus-line germ, the corresponding
actual affine coordinate maps are
\[
 (u,v,t)\longmapsto
       \left(t^2,t^3+\frac{v}{2t^3},\frac{u}{t^2}\right),
 \qquad
 (u,v,s)\longmapsto\left(\frac us,\frac vs,s\right).
                                                        \tag{A.T.21}
\]
Their first generator derivatives are holomorphic
in the respective normalized conormal frames.
The same product differentiation gives cubic pole
bounds six and two, and the fourth derivative bounds
are twelve and four. Thus $\Delta$ has poles bounded
by $12p+4q$ at the two normalization preimages.
The line $L^2(12p+4q)$ is a square and has degree
$2q_0\le0$ by (A.T.19). A nonzero section has no
positive-degree odd divisor, whereas this double
cover has at least four branch points. This is the
last contradiction. All non-lci cases have now
been excluded, and the positive-genus case at the
start gives the asserted elliptic curves and constant
transversal type.

We finally prove the normalization and Euler assertions,
keeping track of any unnecessary blowups in the original
choice of resolution. Factor $\pi$ through the minimal
resolution $\pi_{\min}:Z_{\min}\to T$, with the resulting
sequence of exceptional-curve blowdowns
$q:Z\to Z_{\min}$. Every curve contracted by $q$ is
vertical for $p$, by the stated assumption on $\pi$.
The equality of structure-sheaf direct images for each
blowdown, in local coordinates on $B$, gives a descended
ruling $p_{\min}:Z_{\min}\to B$ with
$p=p_{\min}q$. The image $C_{\min}$ is still a section,
and the pushed-forward effective conductor is
$D_{\min}=2C_{\min}+V_{\min}$ with $V_{\min}$ vertical.
The normalization-canonical comparison gives
$K_{Z_{\min}}+D_{\min}=\pi_{\min}^*\eta^*\omega_S$
as line bundles, and its rational Chern class is zero.

Statement C2d gives a smooth normalization near the
inverse image of the elliptic curve $D$, with normal
line of degree zero on its unramified double cover.
The minimal resolution is an isomorphism there, so
$C_{\min}^2=0$. No other conductor or exceptional
curve meets it. Hence $V_{\min}\cdot C_{\min}=0$.
No ruling $(-1)$-curve can meet $C_{\min}$: one in
$V_{\min}$ is disjoint from it, while one outside
$V_{\min}$ would have
\[
 (K_{Z_{\min}}+2C_{\min}+V_{\min})\cdot E
       =-1+2C_{\min}\cdot E+V_{\min}\cdot E\ge1,
\]
contrary to the vanishing Chern class.

Contract ruling $(-1)$-curves to a minimal ruled bundle.
Every contraction therefore avoids the section. For
one blowdown $\sigma:W\to W'$ let $D'=\sigma_*D_W$,
let $a$ be the coefficient of its exceptional curve,
and let $d$ be the multiplicity of $D'$ at the centre.
The actual divisor comparison is
\[
 K_W+D_W=\sigma^*(K_{W'}+D')+(1+a-d)E.
                                                        \tag{A.T.22}
\]
Its intersection with $E$ gives $a=d-1$, so the
log-canonical line and its zero rational Chern class
descend. The section remains disjoint from the
effective vertical part at every step. On a minimal
ruled bundle every nonzero effective vertical divisor
is a positive sum of fibres and meets the section.
The final vertical part is consequently zero. If
there had been a blowup in the reverse sequence,
its first centre would be away from this section,
and (A.T.22) would give coefficient $-1$ on its
exceptional curve. Its strict transform retains that
negative coefficient under subsequent blowups, a
contradiction to effectiveness. Thus no ruling
blowdowns were needed: $Z_{\min}$ is already a
minimal ruled bundle and $D_{\min}=2C_{\min}$.

All exceptional curves of $\pi_{\min}$ would be
vertical. A minimal ruled bundle has irreducible
fibres of square zero, which cannot form the
negative-definite exceptional divisor of a surface
resolution. Hence $\pi_{\min}$ has no exceptional
curve and is an isomorphism. This proves the assertion
about $T$. Finally $\eta$ is an isomorphism away from
its conductor section $C_T$ and the elliptic image
$D$. Additivity of the topological Euler number on
these actual decompositions gives
\[
 e(S)=e(T)-e(C_T)+e(D)
      =e(\PP^1)e(B)-0+0=0.                       \tag{A.T.23}
\]
The full multiple fibre has the same underlying
topological space and therefore the same Euler number.
\end{proof}

\paragraph{Connected double covers in the two non-lci cases.}
The non-lci calculation also excludes a connected
degree-two conductor cover when its covering curve
is rational, for either of the two image rings in
(A.T.14). In the cusp-plus-line case the cleared
quartic has degree $2q_0\le0$, so it has no branch
divisor at all. In the space-cusp case a rational
double cover has branch degree two. The already
proved inequalities give
\[
                2\le4+2q_0,\qquad 2\le-2q_0+1,
\]
so the integer $q_0$ is $-1$. The cleared quartic
then has degree two and exactly two simple branch
points. Only one point lies above the singular
image germ. The other is a smooth simple branch
point, excluded by the actual coordinate reduction
(A.T.10)--(A.T.12). Thus these two non-lci images
cannot support a connected double cover of any
nonnegative genus in a torus degeneration.

The integer branch budgets, the discriminant square
weights and the coordinate divisibility used in
(A.T.10)--(A.T.12) are checked in
\path{code/E_tacnode_bounds.py}; the non-lci charts,
symbolic generators and cubic/quartic pole orders
are checked in \path{code/E_low_delta_tacnode.py}.
The topological
quotient (A.T.2) targets the actual singular curve
$D/J$; replacing it by its normalization would
lose the free-group obstruction in (A.T.4).


% A.5.19 from part_A6_triple_images.tex

